% !TeX program = pdflatex
\documentclass[review]{fcs}
\usepackage{layout}
\usepackage{subfig}
\usepackage{hyperref}


\title{Frontiers of Computer Science: Instruction for Authors}

% * indicates the authors contributed equally to this work
% + indicates the corresponding author of this work
% 1 or 2 indicates the institutes of the authors of this work
% Use given-name family-name style. Keep hyphens only when they occur in the
% author's own name, e.g., First-name Familyname; do not type family names in all caps.
\author[1,*]{First~Author}
\author[2,*]{Second-name~Author}
\author[1,2,+]{Corresponding~Author}
\address[1]{Higher Education Press, Beijing 100029, China}
\address[2]{Department of Computer Science and Engineering, Beihang University, Beijing 100191, China}

% the email address of the the corresponding author
\corremail{xxx@xxx.com}
\fcssetup{
  received = {month dd, yyyy},
  accepted = {month dd, yyyy},
  article-id = {1},
  first-page = {1},
}

\begin{abstract}
A concise abstract of up to 300 words written in one paragraph, clearly indicating the object and scope of the paper as well as the results achieved, should appear on the first page. References and non-standard or uncommon abbreviations should be avoided as an abstract is often presented separately from the article and must be able to stand alone.
\end{abstract}
\keywords{Large language models; Pre-training; Post-training; Utilization; Evaluation}

\begin{document}

\section{Introduction}

\emph{Frontiers of Computer Science} (FCS) is an international peer-reviewed academic journal sponsored by the Ministry of Education of China. The journal is jointly published by Higher Education Press of China and Springer on a bimonthly basis in English. Online versions of all the contents are available through both \url{http://www.springer.com/journal/11704} and \url{http://journal.hep.com.cn/fcs}.

The journal publishes primarily research and review articles. The submission maybe directly rejected if the field and/or overall quality of the manu\-script are/is not appropriate for our journal. When it is sent for review, normally three reviewers are requested in the relevant study fields.

Manuscripts submitted to this journal must not be under simultaneous consideration by any other publisher and should not have been published elsewhere in a substantially similar form. Not a single part of an article which has been published by FCS may be reproduced or published elsewhere without the written permission of the publisher.

The content of the manuscript must not include information whose disclosure will be prejudicial to the national interest. The authorization license has to be submitted, by the organization the research is sponsored, in case that any information in the manuscript is suspicious to be confidential.

The following categories of papers can be submitted to the journal:

\begin{itemize}
\item RESEARCH ARTICLES describe original investigations relevant to Computer Science.
\item REVIEW ARTICLES give an in-depth over\-view of certain topics or a review of the work of one laboratory or a group of researchers. Review articles usually provide readers with an appreciation of the importance of the work, historical context, a summary of recent developments, and a starting point in the specialist literature.
\item LETTERS are concise articles (no more than 3 printed pages) either discussing topics published in FCS, or reporting original research focused on an outstanding finding or issues of general interest. Different form regular articles, no abstract is required for Letters, and the references should be no more than 10.
\item PERSPECTIVES provide insights, in the form of short articles, on certain topics written by first-class experts in the field.
\item NEWS \& HIGHLIGHTS inform readers about the latest advances in computer science research, as reported in recently published papers or at scientific meetings. News \& Highlights may be peer-reviewed at the editors' discretion.
\item VIEWPOINTS mainly focus on the viewpoints and comments on a specific research topic, provide deep insight into a research problem, or express scientific analysis/views on a computer science topic. Viewpoints may be peer-reviewed at the editors' discretion.
\item COMMENTS can focus on policy, science or other issues related to computer science. The main criteria are that they should be of immediate interest to a broad readership and should be written in an accessible, non-technical style. COMMENTS may be peer-reviewed at the editors' discretion.
\item This journal also publishes special/guest-edited issues. The peer review process for these articles is the same as the peer review process of the journal in general. Additionally, if the guest editor(s) authors an article in their special issue, they will not handle the peer review process.
\end{itemize}

Please kindly note that the extended version of a conference paper that has been published before will not be acceptable in general.

\section{Manuscript preparation and submission requirements}

In this section,  ...

\subsection{Manuscript preparation}

Manuscripts should be in a Word or \LaTeX\ format. The following components are required for a complete regular manu\-script: Title, Author(s), Author affiliation(s), Abstract, Keywords, Nomenclature (if needed), Main text, Acknowledgements, Competing interests, References, Appendices (if need\-ed), Figure captions, Tables. There is no formal limit for the length of a paper, but the editors may recommend condensation when appropriate.

\paragraph{Table}

Large tables should be avoided. If many data are to be presented, an attempt should be made to divide them into two or more tables.
A table should be submitted as three-line tables, that is, there are three horizontal lines: one under the legend, one under the column heads, and one below the body. Vertical lines are generally not used.

\paragraph{Figures}

Figure requirements:
\begin{enumerate}[label=\arabic*)]
\item Numbering and title: number all figures (graphs, charts, photographs, and illustrations) in the order of their citation in the text and cited as, e.g., \cref{fig:1}. Include a title for each figure. Use (a), (b), ... to give titles for subfigures if there are any.
\item Figure quality: should be sharp, noise-free, and of good contrast. All lettering should be large enough to permit legible reduction. The figure quality should meet the requirements: 300dpi for color images; 600dpi for mono images; 72dpi for screen images; 600dpi for line-drawings.
\item Figure borders and size: if a figure has frame lines or boundary boxes, make sure the lines remain visible after reduction. Oversized figures should be scaled down slightly, commonly to about 0.90--0.95\verb|\linewidth| in a single column.
\item Color of figures: The journal will publish all figures in color and free of charge. The color scheme of the figures should be scientifically grounded while maximally embodying artistic qualities, delivering an aesthetically pleasing experience to the audience. Adhere to the color guide as follow.
\begin{center}

\includegraphics[width=.8\linewidth]{color-guide}
\end{center}
\item Figures should be embedded in word or LaTeX processing documents. When accepted, original files should be provided in TIFF, EPS or Corel-Draw file formats.
\end{enumerate}

\cref{fig:1} is one example (example taken from an article published in FCS with some modifications for needed illustrate).

\begin{figure}[!htb]
    \centering
    \includegraphics[width=0.98\linewidth]{fig1.pdf}
    \caption{MRR results and entity frequency grouped by entity in-degree. (a) Result and entity frequency on   NELL23K; (b) result and ...}
\label{fig:1}
\end{figure}

\paragraph{Formulae and equations}

The requirement for formulae and equations are:
\begin{enumerate}[label=\arabic*)]
\item Formulae should be typewritten whenever possible.
\item It is extremely important that all mathematical symbols and letters used are identified and listed, and that the required style of appearance of such symbols is clearly indicated, e.g., bold face, italics, script, outline, etc.( Typically, variables are set in italics.)
\item The words Equation or Equations should appear in full at the beginning of sentences but be abbreviated to Eq. or Eqs. elsewhere.
\item Equations should be located separately from other lines if they are long or complicated (with punctuation mark at the end). Authors are encouraged to place them in two-column format, unless it is necessary to place them in one-column.
\end{enumerate}

For example:
\begin{equation}H=-\frac{1}{2} \Delta \sigma_{\infty}+\sum_{k} \omega_{k} b_{k}^{\dagger} b_{k}+\frac{1}{2} \sum_{k} g_{k}\left(b_{k}+b_{k}\right) \sigma_{z}.\end{equation}

\paragraph{Definition, example, theorem, proof, etc.}

Please do not use italic in these parts except the variables. For example:

\begin{example}
We take the data set 1 ...
\end{example}

\paragraph{Abbreviations}

Expand all abbreviations at their first mention in the text (e.g., architecture analysis and design language (AADL)). Ensure consistency of abbreviations throughout the article.

\paragraph{Footnotes}

Footnotes should only be used if absolutely essential. In most cases it will be possible to incorporate the information in normal text. If used, they should be indicated by superscript Arabic numbers (e.g., 1), 2), ...)\footnote{Footnote information 1}, and kept as short as possible.

\paragraph{Citations and references}

In-text citations must agree with the references in numbering. The references should be presented completely and without mistakes, and should be the original publication. We use the citation-sequence system, i.e., to number the citations according to sequence of their first appearance. Journal names should be spelled out in full. Examples are given in \cite{1,2,3}.

The digital object identifier (DOI) may be used to cite and link to electronic documents, particularly ``Articles in press''.

\subsection{Manuscript submission}

Authors are encouraged to submit their papers electronically via the online submission system (\url{http://mc.manuscriptcentral.com/hepfcs}). The web site guides authors stepwise through the creation and uploading of various files. When you first submit your paper, please upload the PDF file. The manuscript will be finally judged by the editors and the EIC who have the right to accept or reject a paper. When your paper is accepted, you will be required to send all the original source files. All correspondence, including the editor's decision and request for revisions, will be sent by e-mail.

Please note that when submitting the manuscript, all authors' contacts (e.g., e-mail address and affiliations) should be filled in the system. Only the submission is confirmed by all authors can the manu\-script be sent for review.

FCS can also help publish supplementary files (e.g., detailed theoretical proof process, detailed experimental results, etc.) on the journal's website, with no limitation of length.

\section{Copyright transfer}

FCS is a hybrid journal. Authors have Open Access (\textbf{Strongly Recommend}) and Traditional publishing model options to choose how their article is published. Once the choice is made, authors need to sign the corresponding statement. A copy of the statement will be provided with the letter of acceptance of the manuscript. Authors are asked to return by email the signed statement to the editorial office.

Publishing open access offers a number of benefits, including greater reach and readership for your work. We recommend authors to OA. Copyright and license term is CC BY. The article processing charges (APCs) vary by journal---\href{https://www.springernature.com/de/open-research/journals-books/journals}{view the full list}.

\section{Highlights}

Highlights are a short collection of information (e.g., text and graphics), in 3-page PPT, to convey the core findings and provide readers with a quick over\-view of the article. The highlights describe the essence of the research (e.g., results or conclusions) and highlight what is distinctive about it. Highlights will be displayed online (both \url{http://www.springer.com/journal/11704} and \url{http://journal.hep.com.cn/fcs} (Supplementary Material)), but will not appear in the article PDF file or print.

\section{Proofreading and production}

Proofs will be sent to the author and should be returned within 3 days of receipt. Authors should clarify any questions of the proof in a query file. No new materials shall be inserted during the process of proofreading. Please note that authors are urged to check their proofs carefully before returning one all-inclusive e-mail, since subsequent additional corrections will not be possible.

\section{Author enquiries}

For submission inquiries, tracking articles and any information please contact the editorial office. All correspondence for the journal should be sent to \url{FCS@pub.hep.cn}. Please be sure to include the manuscript dispatch number in all correspondences.

% NOTE: The font size of following sections should be 8.5pt.
\fontsize{8.5pt}{10pt}\selectfont

\section*{Acknowledgements}
The ``Acknowledgements'' section is the general term for the list of contributions, credits, and other information included at the end of the text of a manuscript but before the references. Authors should obtain written permission to include the names of individuals in the ``Acknowledgements'' section.

This LaTeX template is formated and maintained by \href{https://www.onlinelatex.com}{OnlineLaTeX}, which is an online LaTeX platform for writing scientific documents such as thesis, papers, resumes, books, reports and letters in LaTeX. Authors are recommened to write their papers on \href{https://www.onlinelatex.com}{OnlineLaTeX} by directly opening this LaTeX template at \url{https://www.onlinelatex.com/models/uppage?tpIds=244}. The authors are required to register an account on \href{https://www.onlinelatex.com}{OnlineLaTeX} first. If you have any questions about the useage of this LaTeX template, please feel free to contact \href{https://www.onlinelatex.com}{OnlineLaTeX} at \url{support@onlinelatex.com}.

\section*{Competing interests}
Conflicts of interest and financial disclosures must be listed in this section. Example of statements to be used when authors have nothing to declare: The authors declare that they have no competing interests or financial conflicts to disclose.

\section*{Appendixes~(if needed)}

\subsection*{Appendix A}

This is Appendix A.

\subsection*{Appendix B}

This is Appendix B.

\bibliographystyle{fcs}
\bibliography{ref}

\begin{biography}{photo}
Please provide each author's biography here with no more than 120 words. The photo (600 dpi, smiling, energetic, front view) can be informal. Our journal prefers to exhibit an encouraging atmosphere. Please use one that best suits our journal.
\end{biography}

\end{document}
